% ЗАКЛЮЧЕНИЕ, без номер.
\section*{ЗАКЛЮЧЕНИЕ}
\label{sec:conclusion}
\addcontentsline{toc}{section}{ЗАКЛЮЧЕНИЕ}

Проектът решава задачата за съхранение на RDF триплети и изпълнение на заявки над тях във вид,
подходящ за вграждане в приложение. Архитектурата разделя синтаксиса, съхранението, анализа,
планирането и изпълнението на слоеве с тесни граници, което позволява всеки да бъде проверяван
поотделно, а изграждането става с компилатор за C++20 и CMake, без външни зависимости.

\textbf{Предимства.} Речниковото кодиране и подредените пермутирани индекси свеждат
съединяването до обхождане на диапазони и премахват низовите сравнения от горещия път. Трите
пермутации покриват всичките осем комбинации от свързани позиции, което е проверено, а не
прието. Точната мощност се получава от самия индекс за $O(\log n)$, така че планировчикът
работи от истински числа.

\textbf{Ограничения.} Индексите се държат изцяло в паметта и се изграждат наново при всяко
отваряне, а трикратното повторение струва 72 от измерените 141 байта на триплет. Поддържаното
подмножество на SPARQL не покрива именувани графи, подзаявки, пътища по свойства, \code{BIND},
\code{VALUES} и \code{HAVING}, а извеждането по RDFS е само в режим на материализация.
Сравнението с утвърдени реализации не е проведено и причината е записана в
Раздел~\ref{sec:results}, а не заобиколена с предположения.

\textbf{Изводи.} Планирането на реда на съединенията си заслужава точно при веригите от три и
повече шаблона, където отношението между непланирания и планирания ред е между 447 и 465 пъти,
докато при един или два шаблона то е в границите на шума; разходът за планиране е под една
стотна от милисекундата, откъдето следва, че планировчик от този вид няма смисъл да се прави
изключваем в производствен код. Стратегията на съединяване се оказа второстепенна, което
противоречи на очакването при проектирането. За разглеждания случай на употреба измерените
величини са в приемливи граници: около 800\,000 до 1\,000\,000 триплета за секунда при
зареждане, около 141 байта на триплет и под три стотни от милисекундата за селективна заявка
над 209\,202 триплета. Заявките от порядъка на десет милисекунди връщат двадесет хиляди
решения, тоест разходът е в резултата, а не в достъпа до данните.

\textbf{По-нататъшна работа.} Отразяване на индексите във файлове; компресия; извеждане по RDFS
по време на заявка; именувани графи; и сравнението с външни реализации, за което работното
натоварване вече е готово.
